\documentclass[11pt,a4paper]{article}
\usepackage[margin=2.5cm]{geometry}

\usepackage[T1]{fontenc}   
\usepackage[utf8]{inputenc} 
\usepackage[english]{babel}

\usepackage{newtxtext,newtxmath}

\usepackage{microtype}

\usepackage{xcolor}
\definecolor{citationred}{RGB}{150,0,0}

\usepackage[
  colorlinks=true,
  citecolor=citationred,
  linkcolor=black,
  urlcolor=black
]{hyperref}

\usepackage[authoryear,round]{natbib}



\usepackage[most]{tcolorbox}  
\tcbuselibrary{skins}         


\newtcolorbox{commentbox}{
  colframe=green,        
  boxrule=0.4pt,         
  sharp corners,         
  enhanced,              
  left=2mm,right=2mm,    
  bottom=2mm,
  top=1.9em,             
  overlay={
    \node[anchor=north west,fill=citationred,text=white,
          font=\bfseries,inner ysep=2pt,inner xsep=4pt]
         at (frame.north west) {Commentaires};
  },
}

\renewcommand{\thefootnote}{\fnsymbol{footnote}}

\title{\textbf{The acceptability of internationally redistributive policies}}
% \thanks{
%     We thank [names] for helpful comments and suggestions. 
%     Any remaining errors are our own.
% }

\author{
    Adrien Fabre\thanks{
        CNRS, CIRED.
        Email: \href{mailto:adrien.fabre@cnrs.fr}
        {adrien.fabre@cnrs.fr}.
    }
    \and
    Lubin Bénéteau\thanks{
        ENS Paris-Saclay. 
        Email: \href{mailto:lubin.beneteau@ens-paris-saclay.fr}
        {lubin.beneteau@ens-paris-saclay.fr}.
    }
}

\date{August 2026}


\begin{document}

\maketitle
\setcounter{footnote}{0}


\section*{Introduction}

A modest reallocation of resources from richer to poorer countries would be enormous relative to the size of low-income economies. Transferring just 1\% of high-income countries' aggregate output to low-income countries (LICs) would be equivalent top double their national income\footnote{According to \href{https://data.worldbank.org/indicator/NY.GDP.PCAP.PP.CD?contextual=default&end=2024&locations=EU-ZG-XD-XM-1W-IN-US-CD-BI-LU-CN&start=2024&view=bar}{World Bank data for 2024}, GDP per capita in high-income countries (HICs) is 28 times higher than in low-income countries (LICs) at Purchasing Power Parity (PPP), and \href{https://data.worldbank.org/indicator/NY.GDP.PCAP.CD?end=2024&locations=EU-ZG-XD-XM-1W-IN-US-CD-BI-LU-CN&start=2024&view=bar}{68 times higher} in nominal terms. Approximately 625 million people live in the 25 low-income countries, defined here as countries with GDP per capita at PPP below \$1,135 per year, compared with 1.42 billion people living in high-income countries. On this basis, 1\% of HICs' GDP represents approximately 60\% of LICs' combined GDP using PPP values and 153\% using nominal values.}. A transfer of this magnitude would be more than sufficient to finance the eradication of extreme poverty \citep{sahoo_what_2025}.

The magnitude of these disparities, together with the potential impact of even limited redistribution, raises a broader question of global justice. In broad terms, global justice concerns the fair and sustainable distribution of socioeconomic and environmental resources both within and across countries, subject to the ecological constraints imposed by planetary boundaries. On this view, reducing extreme poverty and narrowing large international inequalities are necessary but insufficient objectives: distributive improvements cannot be considered sustainable if they are achieved through development paths that intensify climate change and expose already vulnerable populations to further environmental harm. Global justice therefore involves not only the distribution of income and wealth, but also access to public services, economic opportunities, environmental resources, and decision-making power.

International redistribution appears necessary not only to eradicate extreme poverty, a central objective of the universally endorsed Sustainable Development Goals, but also, over the longer run, to achieve convergence in GDP per capita across countries by 2100 \citep{fabre_JDS_2025}. The Global Justice Report estimates that this would require annual transfers from high-income to low-income countries equivalent to approximately 0.8\% of world GDP over the period 2026--2100 \citep[Figure~4.2]{bothe_et_al_2026}. The magnitude of existing global inequality underscores the potential role of such transfers. The richest 1\% of individuals in high-income countries receive more income than the entire bottom half of the global income distribution combined, while the richest 10\% receive more than half of global income \citep{chancel_world_2026}.

These reparative and institutional considerations are complemented by a distinct framework for international redistribution. Within an optimal-tax framework, \citet{kopczuk_limitations_2005} show that a centralized world income tax based on a border-neutral social welfare function would reduce the global Gini coefficient from 0.69 to 0.25, whereas decentralized national systems have almost no effect on global inequality. They conclude that efficiency costs alone cannot explain the limited scale of international redistribution, which is also constrained by free-riding among donor countries.

Moreover, the international distribution of resources cannot be understood solely through differences in current income or gross domestic product. Large flows of value continue to move from the Global South toward the Global North through unequal economic relations, debt servicing, tax-base erosion, and other forms of extraction \citep{hickel_plunder_2021}. Debt distress and austerity can constrain the policy choices available to low-income countries and reproduce asymmetric economic relationships \citep{fischer_return_2023}. The literature also emphasizes the historical origins of contemporary distributive obligations. Colonial institutions contributed to the formation of present economic structures and continue to influence development trajectories \citep{bhambra_colonial_2021}, while slavery and colonial extraction provide distinct grounds for reparative claims \citep{solbiac_britains_2015,darity_cumulative_2022}. Climate change adds a further dimension: because high-income countries have used a disproportionate share of the global carbon budget and, in particular, have already emitted beyond the levels compatible with the temperature limits set by the Paris Agreement, they may bear particular responsibilities toward countries that have contributed less to cumulative emissions but remain highly exposed to their consequences \citep{fanning_compensation_2023,Hickel_and_Zoomkawala_2026}. Taken together, contemporary inequality, historical responsibility, asymmetric economic relations, and climate liability provide several distinct arguments for transferring resources across national borders.

International redistribution should not, however, be conflated with all cross-border resource flows. Migrants' remittances, for example, are private household transfers rather than collectively determined redistributive policies. Officially recorded remittances to low- and middle-income countries reached approximately 0.61\% of world GDP in 2024, exceeding foreign direct investment and Official Development Assistance (ODA) combined \citep{ratha_remittances_2024}. Corporate profit shifting represents a different phenomenon: rather than a deliberate transfer toward poorer countries, it erodes national tax bases and limits governments' ability to finance public expenditure. Around 36\% of multinational profits are estimated to be shifted to tax havens, while associated global corporate-tax revenue losses have been estimated at 0.57\% and 0.74\% of world GDP each year,  although these estimates remain methodologically contested \citep{torslov_missing_2023, cobham_global_2018}. ODA is more directly redistributive and more fully institutionalized, but substantially smaller in aggregate terms. Development Assistance Committee donors provided 0.18\% of world GDP in ODA in 2024 on a grant-equivalent basis, followed by a preliminary decline to 0.15\% of world GDP in 2025 \citep{oecd_final_oda_2025, oecd_preliminary_oda_2026}. 
These comparisons indicate that the politics of international redistribution cannot be reduced to foreign aid alone. It also concerns international taxation, fiscal sovereignty, debt, climate finance, and the institutional rules governing cross-border resource allocation.

This chapter does not seek to resolve the normative question of how much international redistribution justice ultimately requires. It addresses a narrower positive question: to what extent citizens are prepared to accept the policies and institutions through which international redistribution could be implemented? The available literature primarily examines respondents in high-income or net-contributing countries, although several studies provide broader cross-national or globally representative evidence. The chapter therefore focuses principally on the attitudes of citizens  believe that they will bear, part of the costs of international redistribution.

The chapter starts from precisely this puzzle. International redistribution appears normatively important and economically relevant, yet it remains rarely implemented. The argument developed here is a pattern of conditional acceptance. Support tends to be greater when policies are perceived as fair, effective, fiscally credible, progressively financed, and embedded in plausible arrangements for international burden-sharing. Conversely, approval may remain politically weak when citizens doubt that resources will reach intended beneficiaries, perceive foreign expenditure as competing with domestic redistribution, or are held back by the costs involved.

In this chapter, we develop this argument in seven steps. We first examine moral universalism as a psychological and political foundation for extending distributive concern beyond national borders. We then turn to foreign aid, the most familiar public instrument of international redistribution, and studies 
the level of support and how support depends on information, perceived effectiveness, institutional credibility, ideology, and competition with domestic expenditure. In the third section, we consider rule-based forms of redistribution, particularly international taxation of rich individuals funding poverty reduction and public services in low-income countries.
We analyse then moves from instruments to institutions. In the fourth section, we examine whether citizens accept the forms of global governance required to coordinate contributions, constrain free-riding, and implement policies across countries. 

The fifth section turns to international carbon pricing, which would make ordinary citizens in high-income countries bear the costs, and not only the richest people. The sixth section distinguishes collective policy approval from stated readiness to contribute personally, paying particular attention to the gap between respondents’ own propensity to contribute and their beliefs about the propensity of others to do so. 
Finally, the chapter considers why internationally redistributive and climate policies remain politically marginal despite their conditional acceptability. It examines low issue salience, factual misperceptions, second-order beliefs, reference groups, fiscal competition, biased perceptions among political representatives, and weaknesses in political supply.

\section{Deep values: the prevalence and limits of moral universalism}

ISSP surveys indicate that concern for global problems is displayed by a majority. Across 36 countries, 82\% of respondents agree that environmental problems require international agreements that countries must follow, while only 4\% disagree. In a more recent survey covering 29 countries, 78\% respondents consider existing economic differences between rich and poor countries too large, while 5\% who disagree \citep{issp_environment_2010,issp_inequality_2019}. 

Despite such concerns, there is a gap between attitudes toward global inequalities and actual policies. A first attempt to explain it comes from misperception. Indeed, Germans underestimate their global income rank by an average of 15 percentage points. Correcting perceived national rank changes, both national and global, redistributive preferences among left-leaning respondents, whereas correcting perceived global rank has no detectable effect on support for global redistribution. In addition, 90\% of respondents expressed a preference for some degree of global redistribution, although intensity vary across individuals.
By contrast, aid support in US is driven by information on global distribution, because people underestimate their rank by 27 percentage points in their global income rank and overestimate global median income by a factor 10, they tend to increase their support to foreign aid by almost 10 percentage points and donation to international charities by 55\% \citep{nair_misperceptions_2018}.

Literature has long studied the concept of global identity (\citet{rey_2018} for a review). Universalism has several definitions in the literature, although some are extremely close. \citet{enke_moral_2023}, \citet{enke_universalism_2024}, and \citet{cappelen_universalism_2025} define it as the extent to which altruism remains invariant to recipients' identity, group membership, or social and geographic distance. \citet{boyer-kassem_how_2026} retain this conceptual definition but argue that its measurement requires recipients to be economically comparable and to belong to clearly identified groups. A related concept is the moral circle \citep{waytz_ideological_2019,jaeger_toward_2026}.

\citet{bayram_2015} reports that ``78\% of the participants in 57 countries see themselves as citizens of the world'', referring to the 2005--2008 wave of the World Values Survey. In the 2017--2022 wave, however, 84\% of respondents report feeling close or very close to their country, compared with 47\% feeling close to the world\footnote{These results were calculated using the World Values Survey's online analysis tool, available at \url{https://www.worldvaluessurvey.org/WVSOnline.jsp}.}. 

Universalism is reflected not only in stated preferences but also in observed patterns of giving. Using approximately four million donations made through the DonorsChoose platform, worth about \$305 million, \citet{enke_universalism_2024} measure universalism at the U.S. congressional-district level as the extent to which charitable giving remains invariant to geographic and social distance. This district-level measure predicts Democratic vote shares, legislators' roll-call voting, and the moral content of congressional speeches more strongly than district income or education. The individual-level allocation tasks come from a separate paper. \citet{enke_moral_2023} ask 11,000 respondents in Australia, France, Germany, Sweden, and the United States to allocate hypothetical money and trust points between equally rich recipients at different social distances. Their composite measure ranges from zero, when all money and trust are allocated to in-group members, to one, with 0.5 representing equal treatment. Average scores are below this equal-treatment benchmark in every country, ranging from 0.38 in Australia and Sweden to 0.41 in France, with intermediate values of 0.39 in the United States and 0.40 in Germany.\footnote{\citet{boyer-kassem_how_2026} argue that this design does not fully isolate social distance. Equal nominal incomes do not necessarily imply comparable purchasing power, wealth, or relative income positions across countries. Moreover, a ``random stranger'' may share one or more relevant group memberships with the respondent. Differences in allocations may therefore reflect economic disparities or ambiguously defined group boundaries in addition to universalism.}

\citet{prather_2013} controls for inequality and still finds a nationalist bias in the United States, where 84\% of Americans agree that ``taking care of problems at home is more important than giving aid to foreign countries'' \citep{PIPA2001}. Support for a U.S. program to fight hunger is 10 percentage points lower when the hypothetical program operates abroad rather than domestically, and the gap reaches 20 percentage points when the choice concerns cutting an existing program.

Using a similar approach, \citet{enke_moral_2023} show that moral universalism as the extent to which individuals exhibit the same level of altruism and trust toward strangers as toward members of their in-groups, and explains political views better than standard political economy variables such as income, wealth or perception about government efficiency across Australia, Germany, France, the US, and Sweden. For instance, if universalism levels were the same for all, individuals' policy position would not be as much polarized as they are. 
Similarly, using surveys covering 60 countries and 64,000 respondents, \citet{cappelen_universalism_2025} ask respondents to distribute the local-currency equivalent of a hypothetical \$1,000 between two recipients with the same standard of living. Each respondent completes two randomly selected domestic tasks comparing a stranger from their country with either a family member, friend, neighbour, co-religionist, or co-ethnic, followed by a foreign task comparing a domestic stranger with a stranger living elsewhere in the world. The authors define fully universalist respondents as those who divide the money equally in every task: approximately 26\% satisfy this criterion, whereas 17\% allocate, on average across the tasks, no more than 20\% to the stranger. Country averages also vary substantially: the share allocated to the more socially distant recipient is approximately 26\% in China, India, and Israel, compared with 46\% in Ethiopia \citep[Figures 3--5]{cappelen_universalism_2025}.

Universalism can be surprising because it may lead people to sacrifice their own well-being for the sake of distant groups of people. In a lab experiment with U.S. students, \citet{cherry_2017} find that 61.1\% of participants support egalitarian policies even when these reduce their own material payoff. Complementary international experimental evidence from a lab experiment matching participants from Norway and Germany with participants from Uganda and Tanzania in a dictator game shows that participants acted as moral cosmopolitans and did not assign importance to nationality in their distributive choices \citep{cappelen_needs_2013}. Likewise, \citet{leiserowitz_2006} reports that a majority of Americans express greater concern about the consequences of climate change for people worldwide (50\%) or for non-human nature (18\%) than for themselves and their families (12\%) or for the United States (9\%). Consistent with these findings, a 2017 Focus 2030 survey shows that 40\% of French respondents agree that reducing poverty in developing countries should be a priority for the European Union, compared with only 19\% who disagree.

A complementary way to describe this heterogeneity is through the concept of the ``moral circle'' defined as the set of people or other entities considered worthy of moral consideration.
Across seven study samples with U.S. respondents, \citet{waytz_ideological_2019} show that left-leaning individuals consistently report broader moral circles than conservative respondents. They ask their respondents to rate, on a scale from one to five, how strongly they identify with their community, their country, and humanity as a whole. In the full sample, average identification is slightly higher with humanity (3.25) than with the country (3.09) or community (3.06). Among liberal respondents, defined as positions one to three on a seven-point ideology scale, average identification with humanity is 3.39, compared with 2.99 for the country. Among conservatives, defined as positions five to seven, the pattern reverses, with averages of 2.64 for humanity and 3.54 for the country \citep[Figure 3]{waytz_ideological_2019}. 

In two studies in the Netherlands, US, UK and Australia, \citet{jaeger_relative_2023} decompose variation in moral concern into three components. Dutch students evaluate the moral standing of 30 animal targets. Systematic differences between targets explain 40\% of the variation: some animals are consistently considered more worthy of moral concern than others. Differences between judges explain 29\%: some respondents generally assign more moral concern than others. Finally, judge-by-target interactions explain 27\%, meaning that particular respondents evaluate particular targets differently from what their general tendencies and the targets' average ratings would predict \citep[Figures 1]{jaeger_relative_2023}.

We also find this structure when respondents are asked which group they advocate for when they vote. Among 8,000 respondents in five Western countries, \citet{fabre_majority_2025} argue that ``When we ask respondents which group they defend when they vote, 20\% choose ‘sentient beings (humans and animals)’, 22\% choose ‘humans’, 33\% select their ‘fellow citizens’ (or ‘Europeans’), 15\% choose ‘My family and myself’ and the remaining 9\% choose another group (mainly ‘My State or region’ or ‘People sharing my culture or religion’). Notably, a majority of left-wing voters choose humans or sentient beings.''. In a broader survey of 12,000 respondents across eleven 
high-income 
countries, \citet{fabre_public_2026} uses a new question to measure universalism asking respondents the group they advocate when voting and ``45\% choose a universalist  response (either “Humans” or “Sentient beings (humans and animals)”), while 32\% opt for their fellow citizens''. The universalist share reaches 50\% in Europe, 57\% in Saudi Arabia, and 59\% among left-wing respondents \citep{fabre_public_2026}. Universalistic concern is widespread, despite median voters favoring their fellow citizens in the U.S. and Japan.

When respondents are asked why ``high-income countries should support low-income countries'', they first invoke a moral duty, on average at 54\% at the exception of France (43\%) and Japan (45\%), while, on average, 38\% of respondents suggest the ``Long-term interest of HICs to help LICs''. Finally, only 16\% of respondents advocate non of these arguments, which puts an end to self-interest or guilt bias \citep{fabre_public_2026}.

\section{Foreign aid: broad but conditional support for increased spending}

Public attitudes toward foreign aid in high-income countries have been studied for several decades. This literature shows that citizens commonly overestimate the amount their country spends on foreign aid, even though majorities generally support maintaining or increasing aid, particularly when it is perceived as effective. In the United States, 83\% of respondents supported a multilateral initiative to halve world hunger \citep{PIPA2001}. Across 20 countries, an average of 81\% agreed that developed countries ``have a moral responsibility to work to reduce hunger and severe poverty in poor countries'' \citep{PIPA2008}. Moreover, in each of seven OECD countries surveyed, at least 75\% of respondents were willing to contribute financially to a program designed to halve world hunger, for an estimated cost of, for example, \$50 per year for an American respondent.

\citet{kaufmann_foreign_2019} document that respondents in each of the 24 countries they study overestimate foreign aid, by a factor of seven on average. In most countries, the level of aid people would like their government to provide exceeds the level they believe is currently provided.\footnote{\citet{kaufmann_foreign_2012} provide the most reliable estimates of desired aid because, as noted by \citet{hudson_mile_2012}, earlier studies generally failed to account for misperceptions of actual aid levels.} Respondents in the top income quintile support aid levels that are 0.13 percentage points of GNI lower than those preferred by the bottom 40\%. Combining a theoretical framework with data on lobbying and observed aid allocations, while controlling for citizens' preferences in terms of desired aid, the authors conclude that actual aid falls short of desired aid because wealthier individuals exert disproportionate political influence toward their economic interests. The United States stands out as an exception. Desired aid is similar to the average across other countries, at around 3\% of GNI, but respondents believe that aid is already about twice as large as the level they would ideally support. Perceptions of current aid alone explain more than 15\% of the variation in desired aid. From one country to another, the ratio of median desired aid to actual aid ranges almost from parity to a ratio of more than 9 to 1, revealing what the authors describe as a substantial ``democratic aid deficit'' in several donor countries.

The United States is an important exception in the cross-country evidence. In the data analyzed by \citet{kaufmann_foreign_2019}, Americans preferred foreign aid equivalent to 3.19\% of GNI, close to the cross-country average, but believed that their country was already spending 7.52\%, compared with an actual level of only 0.18\%. This particularly large misperception is consistent with earlier U.S. evidence. \citet{gilens_political_2001} estimates that, among respondents with high levels of general political knowledge, informing them that foreign aid accounted for less than 1\% of federal spending reduced the predicted share favoring aid cuts from 58.6\% to 42.0\%. Likewise, \citet{nair_misperceptions_2018} attribute weak support for aid among US-Americans to widespread misperceptions of their rank in the global income distribution: respondents underestimate their rank by 27 percentiles on average while they overestimate the global median income by a factor of ten. Such misperceptions may explain why, in the 2000-2004 waves of the GSS, more than 60\% of Americans believed that the government was spending too much on foreign aid \citep{okten_preferences_2007}.

Using a survey experiments in the United States, the United Kingdom, France, Germany and Spain, \citet{fabre_majority_2025} show that, on average across these countries,
respondents greatly overestimate current foreign aid (e.g., 4.7\% against an actual 0.4\% of public spending in the U.S.), but still prefer substantially higher aid budgets than those currently observed, allocating on average 1.8\% after receiving information on the true amount and 4.0\% without it \citet[Figure 8]{fabre_majority_2025_si}.

Respondents were also asked how much they believed their country spent on foreign aid. The median respondent overestimated the actual amount. In the United States, the median estimate ranged from 1.8\% to 2.6\%, compared with an actual level of 0.4\%. In the pooled Western European sample, respondents believed that foreign aid represented 2.9\% of public spending, compared with an actual level of 1.1\%. Their mean preferred allocation was 3.9\% without information and 2.7\% after information was provided \citep[Figure S20]{fabre_majority_2025_si}.

Respondents were then asked about their attitudes toward their country's evolution of foreign aid. 60\% of Americans are in favor for an increase, 43\% of them wants conditions attached to this, while the corresponding figures for the pooled Western European sample are 64\% and 39\% to 51\% \citep[Figure 6]{fabre_majority_2025_si}. In addition, respondents were asked about the government spending they would allocate to foreign aid, with a random subgroup first informed of the actual amount. Without this information, the median preferred interval was 1.8-2.6\% in the United States, France, and Germany, compared with 1.1-1.7\% in Spain and the United Kingdom. Informing respondents reduced the U.S. median to 0.3-0.5\%, while the French and German medians remained at 1.8-2.6\% \citep[Figure S20-S21]{fabre_majority_2025_si}. Respondents, for about 47\% in U.S., 59\% in Europe and even 75\% in France ``Preferred foreign aid is higher than current'' \citep[Figure S23]{fabre_majority_2025_si}. Respondents who preferred higher foreign aid were asked how the increase should be financed. ``Higher taxes on the wealthiest'' were by far the most frequently selected option, chosen by 68\% of Americans and 64\% of Europeans, and by as many as 82\% in Spain and 85\% in the United Kingdom. By contrast, relatively few, between 1\% and 8\% respondents favored cuts in education \citep[Figure S24]{fabre_majority_2025_si}. Respondents who preferred lower aid were asked how the budgetary savings should be used instead. Healthcare services spending was generally prioritized: 57\% of Europeans favored it while it represents 40\% of Americans \citep[Figure S25]{fabre_majority_2025_si}.

Respondents who are in favor of an increase in foreign aid only under certain conditions were asked: ``Under which conditions should foreign aid be increased?''. \citet{fabre_majority_2025} find that the most important condition is ensuring that ``the aid reaches people in need and money is not diverted'', supported by 77\% of Europeans and 68\% of Americans. A second widely supported condition is ``that recipient countries comply with climate targets and human rights'' (72\% in Europe and 61\% in the United States). Around half of respondents also require that ``other high-income countries also increase their foreign aid'' (51\% in Europe and 45\% in the U.S.), whereas conditioning aid on ``increased taxes on millionaires'' receives more limited support (38\% and 36\%, respectively) \citep[Figure 7]{fabre_majority_2025}.

Reviewing the literature, \citet{hudson_mile_2012} conclude that support for foreign aid is largely motivated by an intrinsic concern for poverty, while also documenting considerable skepticism about aid effectiveness. According to \citet{dfid_aid_2009} and \citet{PIPA2001}, 47\% of British respondents believe that aid is wasted, primarily because of corruption, while Americans estimate that less than one-quarter of aid reaches people in need and that more than half is captured by corrupt government officials. These beliefs do not, however, rule out a utilitarian motivation for supporting aid: even if only a fraction of the funds reaches the intended recipients, the expected welfare gains from transferring resources to substantially poorer populations may remain large \citep{kopczuk_limitations_2005}.

Consistent with \citet{henson_public_2010}, \citet{bauhr_corruption_2013} find that perceived corruption in recipient countries lowers support for aid. However, this relationship is mitigated by what they describe as the aid-corruption paradox: countries where corruption is major are often those with the highest need for assistance. \citet{bodenstein_who_2017} further show that right-leaning Europeans and respondents who perceive high levels of corruption are more likely to support conditioning EU development assistance on recipient countries complying with standards of democracy, human rights, and good governance.

Respondents who preferred reducing or maintaining foreign aid were asked: ``Why should foreign aid should not be increased?'' The most common reason is ``Charity begins at home: there is already a lot to do to support the [country] people in need'', selected by 63\% of Europeans and Americans who refuse to increase aid. The second most frequent answer is ``Aid is not effective as most of it is diverted'', chosen by slightly half of these respondents %53\% of Europeans and 40\% of Americans 
\citep[Figures 7, S25]{fabre_majority_2025}.

Using the 2002 Gallup survey and the 2006 World Values Survey, \citet{bayram_aiding_2017} and \citet{paxton_individual_2012}, consistent with the U.S. evidence of \citet{heinrich_voters_2018}, identify trust, left-wing political orientation, political interest, and female gender as the main individual-level correlates of support for higher aid.

Using Eurobarometer data from 15 OECD donor countries, \citet{cho_ideology_2024} shows that the association between ideology and the intensity of moral support for foreign aid is stronger among socioeconomically advantaged respondents than among disadvantaged respondents. At the far-left end of the ideological scale, the predicted probability of answering ``totally agree'' that tackling poverty abroad is a moral obligation is approximately 63\% among advantaged respondents, compared with 44\% among disadvantaged respondents. Conversely, the probability of answering ``tend to agree'' is around 31\% among advantaged respondents and 42\% among disadvantaged respondents. Thus, among the advantaged left, the highest level of support clearly predominates, whereas among the disadvantaged left, ``totally agree'' and ``tend to agree'' are nearly equally prevalent. As respondents move to the right, ``tend to agree'' becomes the most prevalent response among advantaged respondents, while its predicted probability remains relatively stable across the ideological scale among disadvantaged respondents. Ideology therefore more strongly differentiates the intensity of support among advantaged respondents, whereas disadvantaged respondents tend to express more moderate support across ideological positions \citep[Figure 1]{cho_ideology_2024}.

Eurobarometer surveys reveal broad acceptance of foreign aid across the European Union. In 2015, in the 28 EU Member States, approximately 64\% supported honouring the EU's commitment to increase aid or increasing it beyond the promised amount \citep{european_commission_2015_eurobarometer}. 
In 2019, 73\% thought that the EU and its Member States should either spend more on financial assistance to partner countries (30\%) or maintain spending at its current level (43\%). Of the remaining respondents, 17\% favored lower spending and 10\% answered ``Don't know'' \citep{european_commission_2019_development}.
In 2020, a survey across the 27 Member States found that 91\% of respondents considered EU funding for humanitarian aid important \citep{european_commission_2021_humanitarian}.
In 2022, in all 27 Member States, 80\% stated poverty reduction in developing countries should be a priority of the EU, and 67\% considered it should be a priority for their national government \citep{european_commission_2022_development}. 

\section{International taxation: strong support for taxing the rich to reduce global poverty}

\citet{fabre_public_2026} find in surveyed countries that all respondents accept\footnote{There is a distinction between \textit{support} and \textit{acceptance}. The term \textit{support} refers to the absolute proportion of \textit{Somewhat agree} or \textit{Strongly agree} responses on Likert scales, whereas the term \textit{acceptance} refers to relative support, that is, the proportion of support among responses non-\textit{indifferent}.} in majority global redistributive policies tested. The 2\% tax on billionaire wealth is the most supported policy totalling 81\% acceptance. Other proposals as debt reliefs for vulnerable countries -- where developed countries contribute at 0.7\% of the GDP in foreign aid --, expansion of the UN Security Council, or even, the Bridgetown initiative (an sustainable increase of investments at low interests rates in low income countries), gather overall, 70\% of acceptance. 

\citet{fabre_public_2026} tests support for ten plausible policies currently debated in international organizations. Figure \ref{fig:plausible_policies} reports their acceptance\footnote{There is a distinction between \textit{support} and \textit{acceptance}. The term \textit{support} refers to the absolute proportion of \textit{Somewhat agree} or \textit{Strongly agree} responses on Likert scales, whereas the term \textit{acceptance} refers to relative support, that is, the proportion of support among responses non-\textit{indifferent}.} across the surveyed countries. In the pooled sample, every proposal receives majority acceptance, ranging from 53\% for the international levy on aviation emissions to 81\% for the 2\% minimum tax on billionaires' wealth.

\begin{figure}[htbp]
    \centering
    \includegraphics[width=\textwidth]{chapter/images/figure_10_fabre_26.png}
    \caption{Acceptance of plausible global redistribution policies. \textit{Source:} \citet[Figure 10]{fabre_public_2026}.}
    \label{fig:plausible_policies}
\end{figure}

The ``plausible'' label comes from the international debates upon these policies\footnote{As documented in \citet[Appendix C.1]{fabre_public_2026}, where polices as been discussed during summit as G20, during COPs or at ONU and OCDE}. As mentioned above, almost every policy proposal gather a majority of acceptance. Nevertheless, Japanese respondents acceptance on aviation taxation is at 46\% and it could be explained by the salient cost displayed through the 30\% increase in flight prices. For all plausible policies, left-leaning Europeans garner the highest acceptance rates (from 71\% for aviation taxation to 94\% for the 2\% billionaire taxation) followed by the U.S. Harris voters respondents (respectively, 67\% to 93\%). On the opposite, U.S. respondents who voted for Trump and far right voters in Europe display a reduced support compare to the overall mean (e.g., 34\% of Trump voters support aviation tax while 56\% support 2\% billionaire tax, the corresponding numbers are 40\% and 74\% for far-right Europeans voters) \citep[Figure S71]{fabre_public_2026}.

To complete these arguments: acceptance of plausible policies, \citet{cappelen_universalism_2025} ask respondents is the national government provide foreign aid in order to reduce ``economic differences between the rich and the poor'' in selected countries rather ``national government should focus on helping the poor in [country], rather than the poor elsewhere in the world''. Respondents generally prioritize the domestic poor: after responses are reverse-coded so that higher values indicate greater priority for the global poor, the mean lies around 1.5 on a four-point scale \citep[Figure 7B]{cappelen_universalism_2025}.
In the first experiment authors ask respondents to split hypothetical \$1000 between a in-group member (family, friend, neighbor, co-religionist, or co-ethnic) or a stranger in their country. The second experiment have the same method but oppose a stranger from her country to a stranger elsewhere in the world, respondents are explicitly told that the two recipients have the same standard of living. On average, respondents allocate 38.1\% to domestic strangers and 37.5\% to foreign strangers while 27\% divide the money equally in all decisions, whereas 6\% consistently allocate everything to the in-group member. Moving from complete in-group favoritism to equal treatment is associated with a 0.09 point increase in support for reducing domestic inequality and a 0.17 point increase in support for prioritizing the global over the domestic poor, both measured on four-point scales. These estimates are correlations rather than causal effects \cite[Table 1]{cappelen_universalism_2025}.

\citet{fabre_majority_2025} use a donation experiment in the United States and Europe to test support for international redistribution policies. Each respondent has to split a potential \$100 lottery prize between themselves and people living in poverty either in their home country or in Africa, with the recipient's location randomly assigned. In Europe, respondents allocate similar amounts to poor recipients in their own country and in Africa. In the United States, however, respondents give approximately \$3 more to poor recipients in their own country than to those in Africa. The estimated difference is \$2.509 and is statistically significant, whereas the corresponding difference in Europe is not statistically significant \citep[Figure S13 \& Table ED2]{fabre_majority_2025_si}.

In another allocation task, respondents have to allocate 100 points among six randomly selected policies. An equal allocation would therefore give approximately 16.7 points to each policy. On average, respondents in the U.S. and in Europe respectively allocated 21 and 20 points to a ``Global tax on millionaires'', placing it above this benchmark. By contrast, the policy ``Doubling foreign aid'' gathered 11 points, on average, in Europe and 9 in the U.S., placing it below the benchmark \citep[Figure S30]{fabre_majority_2025_si}.

The authors then ask respondents how many policies they support from a randomly assigned list. In the most comprehensive list, respondents evaluate four policies---the Global Climate Scheme, the National Redistribution Scheme, a coal exit, and ``marriage only for opposite-sex couples''---and can therefore report supporting between zero and four policies. The midpoint of this scale is two policies. The mean number of supported policies exceeds this midpoint in all four European countries: respondents support an average of 2.8 policies in France, 2.7 in Spain, 2.6 in the United Kingdom, and 2.2 in Germany. In the United States, the average is exactly at the midpoint, at 2.0 policies \citep[Figure S7]{fabre_majority_2025_si}.

In another allocation task, respondents have to allocate 100 points to different among six randomly selected policies. On average, respondents in the U.S. and in Europe respectively allocated 20 and 21 to a ``Global tax on millionaires''. While the policy ``Doubling foreign aid'' gathered 11 points, on average, in Europe and 9 in the U.S. \citep[Figure S30]{fabre_majority_2025_si}. 

\citet{fabre_majority_2025} asks respondents what share of the revenues from a worldwide tax on household wealth above \$5 million should be allocated to low-income countries. This share would be used to ``finance infrastructure and public services in low-income countries''. The median response is 30\% of the tax revenue across the five countries \citep[Figure S16 and ED4]{fabre_majority_2025_si}.

\citet{fabre_public_2026} conducts a revenue allocation task including five spending items---four of them concern domestic issues: the deficit, income tax, welfare programs, and education, while one concerns healthcare, education, and renewable energy in low-income countries. Respondents are asked how much of the revenue from a hypothetical wealth tax, applying a rate of 2\% to wealth above \$5 million, they would allocate to each item, after being informed of the revenue that the tax would generate in their country and in all low-income countries combined. Fabre finds, on average, that 17\% of the tax revenue would be allocated to global issues. The countries with the highest percentages are Spain and Saudi Arabia, both at 21\%, while the United States is at 17\% and Japan at 14\%.

\citet{fabre_public_2026} examines support for ``a global tax on top incomes to finance poverty reduction in LICs''. Respondents are randomly presented with either a tax targeting the global top 1\% or a more progressive schedule covering the global top 3\%. Under the first proposal, after-tax income above \$120,000 per year at purchasing power parity is subject to an additional rate of 15\%. The second applies marginal rates of 15\%, 30\%, and 45\% above annual income thresholds of \$80,000, \$120,000, and \$1 million, respectively. The revenues would finance transfers associated with monthly poverty thresholds of \$250 under the first proposal and \$400 under the second, corresponding to approximately 2\% and 5\% of world nominal income. Absolute support for the tax on the global top 1\% reaches 56\% overall, 61\% in Europe, and 51\% in the United States \citep[Figure S40]{fabre_public_2026}. Among respondents who would personally pay either tax, 53\% accept the proposal overall; acceptance reaches 61\% in Europe, 32\% in Switzerland, and 79\% in Saudi Arabia. These affected respondents represent 6\% of the pooled sample, including 4\% in Europe, 11\% in Switzerland, and 21\% in Saudi Arabia \citep[Figure S41]{fabre_public_2026}.

\citet{fabre_public_2026} then asked respondents to construct their preferred global distribution of after-tax income by changing the shares of winners and losers and the redistribution from richer to poorer individuals. Respondents see both the current and proposed distributions: ``The current distribution is in red, and your custom one is in green'' and could preserve the status quo. The sliders' initial positions are randomly set to limit, and respondents are allowed to skip the task. The interactive exercise is available at \url{http://preferences-pol.fr/custom_global_redistr.html}, with an explanatory video at \url{https://youtu.be/gSfsQwczT7w}. Overall, 56\% of respondents declare themselves satisfied with their redistribution, whereas 43\% skip the task \citep[Figure S51]{fabre_public_2026}. Among respondents, the median proposal made 49\% of the world population better off and 18\% worse off, with a redistribution intensity of 5 out of 10 \citep[Figure S49]{fabre_public_2026}. This proposal would transfer 5.4\% of world income from richer to poorer individuals and raise the minimum monthly income to \$287 at purchasing power parity \citep[Figure S49]{fabre_public_2026}. Almost half of respondents (48\%) select a redistribution that would reduce their own income, compared with 9\% who would personally gain, while 10\% retain the current distribution. The average preferred redistribution would likewise transfer 5.4\% of world income, in this case from the top 27\% to the bottom 73\%, and establish a minimum monthly income of \$247 \citep[Figure S50]{fabre_public_2026}. The average and median proposals would reduce the global Gini coefficient from 0.67 to 0.59 and 0.58, respectively. At the individual level, the largest group of respondents (41\%) selected transfers representing between 4\% and 5\% of world income, while 55\% chose a redistribution producing a minimum monthly income between \$150 and \$350 \citep[Figure S56-57]{fabre_public_2026}. At the top of the distribution, the average proposal could be implemented through additional marginal tax rates of 7\% above \$25,000 and 16\% above \$40,000 per year at purchasing power parity \citep[Figure S69]{fabre_public_2026}.

Fabre finds original results from an online survey of 1,000 respondents representative of the French population conducted in 2026, showing that 41\% of respondents support a tax on the top 5\%, while 36\% oppose it, support decreases to 37\% for a tax on the top 8\%, while opposition increases to 39\%.

\section{Global governance: public support for democratic decision-making on global issues}

%===============================================================================================
% - International governance needed to achieve int'l redistribution: are people ready?
%===============================================================================================

International redistribution requires an institution capable of setting common rules, mobilizing resources and implementing collective decisions across countries. The key question is therefore whether citizens are willing to grant international institutions enough authority to act beyond national borders, including for redistributive purposes. 

%===============================================================================================
% - works of Farsan Ghassim
%===============================================================================================
%Fabre 25 NHB, p 1584 
\citet{ghassim_who_2024} find that support for world government strongly depends on its institutional design. Their survey covers 17 countries: Argentina, Australia, Canada, China, Colombia, Egypt, France, Hungary, India, Indonesia, Kenya, Russia, South Korea, Spain, Turkey, the United Kingdom and the United States. Using population weights, support reaches 58\% for an unspecified world government, 71\% for a democratic one, 69\% for one restricted to global issues, and 72\% for an institution combining both characteristics\footnote{Responses were recorded on a six-point scale: ``strongly oppose'', ``oppose'', ``somewhat oppose'', ``somewhat support'', ``support'' and ``strongly support''. There was no neutral midpoint. Here, ``support'' includes the answers ``somewhat support'', ``support'' and ``strongly support''.} \citep[Figure 2]{ghassim_who_2024}.

The authors then focus on the fully specified proposal: a democratic world government empowered to address ``climate change, world poverty, and international peace'', while national governments retain authority over non-global issues. With population weights, support reaches 72\% in the full sample, 75\% in partly free or non-free countries but 56\% in free countries. More generally, support is higher in more populous, less developed, less free and less powerful countries \citep[Figures 2 \& 5]{ghassim_who_2024}.

The fully specified proposal receives majority support in 16 of the 17 surveyed countries. In these 16 countries, support ranges from 56\% to 82\%, indicating that approval is not confined to a particular region, income group or political regime. The United States constitutes the only exception, with 55\% of respondents opposing the proposal \citep[Figure 6]{ghassim_who_2024}. Overall, these results document broad cross-country support for a democratic global institution whose authority is limited to transnational issues.

Support for global democracy is sufficiently salient to affect stated voting intentions. Across Brazil, Germany, Japan, the United Kingdom and the United States, between 55\% and 74\% of respondents support a directly elected global parliament \citep[Figure 7]{ghassim_global_democracy_2020}. In Brazil and Germany, informing respondents about parties' supposed positions on global democracy shifts aggregate voting intentions from parties described as opposing the proposal toward those described as supporting it by 8 and 12 percentage points, respectively \citep[Figures 144--145]{ghassim_global_democracy_2020}. In Germany, when respondents are told that only the Greens and Die Linke support global democracy, these parties gain 9 and 3 percentage points, respectively, while the SPD and the CDU-CSU each lose 6 percentage points \citep[Figure 145]{ghassim_global_democracy_2020}.

\citet{ghassim_public_2022} ask respondents in Argentina, China, India, Russia, Spain and the United States which institutional features they prefer for a stronger and more representative UN. Making UN decisions binding on every country in areas such as security, environmental and economic matters receives a marginal mean of approximately 0.53\footnote{A marginal mean is the average probability that a profile containing a given institutional feature is selected.}. 
By contrast, limiting their binding force to countries that ``voluntarily accept'' them is the least preferred option, with a marginal mean of about 0.46. The option of having no enforcement mechanism also receives a relatively low marginal mean of 0.48. Collective enforcement by member states and direct enforcement by a strengthened UN reach approximately 0.51 and 0.52, respectively, although the difference between these two options is not statistically significant. Finally, creating a second chamber composed of representatives directly elected by citizens receives a marginal mean of about 0.53, while assigning all important decisions to a body in which every member country is represented reaches approximately 0.50 \citep[Figure 1]{ghassim_public_2022}.

\begin{figure}[htbp]
    \centering
    \includegraphics[width=\linewidth]{chapter/images/intl_policy_en.pdf}
    \caption{Evaluations of international policy proposals among respondents in France in 2026.}
    \label{fig:intl_policy}
\end{figure}

\begin{figure}[htbp]
    \centering
    \includegraphics[width=\linewidth]{chapter/images/intl_governance_en.pdf}
    \caption{Preferred mechanisms for making decisions on global issues among respondents in France in 2026.}
    \label{fig:intl_governance}
\end{figure}

Fabre reports original evidence from an online survey conducted in 2026 among 1,000 respondents representative of the French population. Overall, 60\% favor involving citizens worldwide in decisions on global issues. Among the proposed mechanisms, 75\% evaluate direct worldwide referendums favorably, while 70\% support representative surveys of global public opinion. Citizen assemblies selected by lot, whose decisions would subsequently be implemented by national parliaments, receive 61\% favorable evaluations. By comparison, 54\% favor a directly elected world parliament responsible for most decisions \url{https://github.com/bixiou/budget_26/blob/main/figures/intl_policy_en.pdf} \url{https://github.com/bixiou/budget_26/blob/main/figures/intl_governance_en.pdf}.


\section{International climate policy: majority acceptance despite personal costs}
%===============================================================================================
% - burden-sharing: every country should participate, burden should be equitably shared
%===============================================================================================

International redistributive mechanisms, including some climate-policy schemes, must determine how countries share the costs. Public acceptance is likely to be higher when a broad coalition participates and contributions reflect each country's capacity.

Survey experiments on climate burden-sharing confirm the importance of institutional design. Support increases when many countries participate and share the costs rather than leaving most of the burden to the respondent's country. Respondents nevertheless remain willing to accept substantial domestic contributions. \citet{bechtel_mass_2013} conducted survey experiments in 2012 with nationally representative samples from France, Germany, the United Kingdom, and the United States. Each respondent completed four comparisons between two randomly generated international climate agreements. For each comparison, respondents selected their preferred agreement and reported how likely they were to vote for or against their country's participation in a referendum. The authors estimate that moving from an agreement resembling the one then under discussion to the design predicted to maximize support would raise the share voting in favor from 42\% to 60\% in France, from 37\% to 60\% in Germany, from 36\% to 51\% in the United Kingdom, and from 29\% to 47\% in the United States \citep[Table S3]{bechtel_mass_2013}.


The agreements varied in their monthly household cost, the distribution of costs across countries, the number of participating countries, the share of global emissions covered, the monitoring institution, and the sanctions for non-compliance. These design features strongly affected support. Increasing costs from 0.5\% to 1\% of national GDP reduced the probability of selecting an agreement by approximately 10 percentage points. By contrast, increasing participation from 20 to 160 countries out of 192 raised it by about 15 points \citep[Table 1 \& Figure 2]{bechtel_mass_2013}.

Focusing on North--South climate finance, \citet{gampfer_obtaining_2014} conducted survey experiments with nationally representative samples in the United States and Germany between October and November 2013. Respondents considered six pairs of hypothetical climate-finance mechanisms and indicated, for each pair, ``which of the two proposals they preferred''. The mechanisms varied in recipient countries' climate damages, income, emissions, and government efficiency; the actors controlling allocation decisions; the use of funds for mitigation, adaptation, or both; and the distribution of contributions between the respondent's country and other donor countries. Donor burden-sharing produced by far the largest effect. Relative to a mechanism in which other countries provided 10\% of the funds and the respondent's country provided 90\%, reversing these shares increased the probability of selection by 28 percentage points in the United States and 23 points in Germany. Respondents nevertheless accepted substantial domestic contributions: most of the increase in support occurred when the share paid by other countries rose from 10\% to 30\% or 50\%. 
When decisions about which projects would receive funding were made jointly by contributing and recipient countries rather than by recipient countries alone, the probability of selecting the mechanism increased by 6 percentage points in the United States and 9 points in Germany. Respondents also preferred mechanisms that financed both mitigation and adaptation to those that financed adaptation alone \citep[Table 1 \& Figure 2]{gampfer_obtaining_2014}.

%===============================================================================================
% - Global Climate Scheme
%===============================================================================================

\citet{fabre_public_2026} examines public acceptance of a Global Climate Scheme (GCS). The description shown to respondents began: ``In 2015, all countries agreed to contain global warming `well below +2\textdegree{}C'.'' The proposed scheme would impose a global ceiling on greenhouse-gas emissions and require polluting firms to purchase permits within that limit. By incorporating the cost of emissions into fossil-fuel prices, the mechanism was expected to reduce fossil-fuel consumption and, consequently, greenhouse-gas emissions. Permit revenues would finance equal monthly payments to adults worldwide, reflecting the stated principle that ``each human has an equal right to pollute.'' Respondents were informed that an average transfer of \$35 per month would lift approximately 600 million people earning less than \$2 per day out of extreme poverty. They were also shown country-specific estimates of the net monthly loss borne by a typical resident and of the corresponding transfer: \$90 and \$35 in the United States, \texteuro 15 and \texteuro 20 in France, \texteuro 45 and \texteuro 20 in Germany, CHF 15 and CHF 20 in Switzerland, and RUB 2,500 and RUB 3,000 in Russia. In the United States, the projected 2\% increase in consumer prices implied a net monthly loss of \$90 after subtracting the \$35 transfer \citep[Table S2]{fabre_public_2026}.

%LB: il faudra juger de l'utilité de ce paragraphe --> la Figure S30 suit ce paragraphe.
The GCS is accepted by 55\% of respondents in the pooled sample. Acceptance is lowest in the United States and Russia, at 49\%. It reaches 61\% in Europe as a whole, including 65\% in France, 53\% in Germany, and 58\% in the United Kingdom. Among the eleven surveyed countries, Saudi Arabia displays the highest acceptance rate, at 85\% \citep[Figures 6 \& S30]{fabre_public_2026}.


%===============================================================================================
% - International Climate Schemes
%===============================================================================================

The International Climate Scheme (ICS) applies the same cap-and-trade mechanism as the GCS, but within different coalitions of countries. Respondents are randomly assigned to one of four variants and shown a map displaying ``a possible set of countries'' that would participate in the scheme, together with their share of global emissions. In the \textit{Low} variant, the coalition comprises the Global South and the European Union, but excludes China, covering 33\% of global emissions. The \textit{Mid} variant replaces the EU with China and brings together China and the Global South, which account for 56\% of global emissions. The \textit{High} variant adds the EU, the United Kingdom, Japan, Canada, South Korea, and other high-income countries to the \textit{Mid} coalition, raising coverage to 72\% of global emissions. Finally, the \textit{High color} variant retains the same coalition as the \textit{High} variant but uses a colored map to display the net financial gain or cost of the scheme for each country. Respondents assigned to this variant were told that ``a provision would prevent the Global Climate Scheme from harming low- and middle-income countries'', so countries such as China ``would neither win nor lose financially'' \citep[Questions 31--35 and Figure 5]{fabre_public_2026}. Wider participation increases acceptance only modestly: support rises from 65\% in the \textit{Low} variant to 68\% in the \textit{High} variant. By contrast, explicitly displaying national gains and losses reduces support by 6.6 percentage points relative to the \textit{High} variant, suggesting that the salience of distributive consequences matters more than the precise size of the coalition \citep[Figure 6]{fabre_public_2026}.

All four ICS variants attract majority support in every country surveyed \citep[Figures 6 \& S30]{fabre_public_2026}. Acceptance of the \textit{Low} coverage variant is similar to that of the National Climate Scheme, although the ICS is more costly for the average respondent. As Fabre notes, this comparison ``suggests that the average respondent is willing to pay the ICS's higher cost'' in exchange for guaranteed poverty alleviation and decarbonization in the Global South. The ICS also receives greater acceptance than the GCS, but this may reflect the greater credibility and precision of a proposal that lists participating countries, its visual presentation, or the lower salience of individual costs. Nevertheless, ``the data does not allow testing these different hypotheses'' \citep{fabre_public_2026}.

\begin{figure}[htbp]
    \centering
    \includegraphics[width=\textwidth]{chapter/images/figure_S30_fabre_26.png}
    \caption{Support for the National, Global, and International Climate Schemes by country \citep[Figure S30]{fabre_public_2026}.}
    \label{fig:ics_support_country}
\end{figure}


%===============================================================================================
% - climate policy supported if it's perceived as fair, environmentally effective and in one's interest, suggesting majority support as soon as two out of three elements are there (Douenne & Fabre 2022); if int'l climate policy seen as fair and effective, logical that it's supported even if against one's interest
%===============================================================================================

%LB: paragraphe raccourci, j'ai quand même gardé deux chiffres.
\citet{douenne_yellow_2022} show that acceptance of climate policy depends strongly on three perceived features: whether the policy is fair, environmentally effective, and in the respondent's own financial interest. In their estimates, believing that one would not lose financially increases acceptance by 53 percentage points, while believing that the policy is environmentally effective raises approval by 42 points. The authors therefore suggest that convincing citizens on more than one of these dimensions would likely generate large majority support, although they cannot causally identify their combined effect. Accordingly, an international climate policy may remain majoritarian even when it imposes a personal cost, provided that it is regarded as both fair and environmentally effective \citep[Tables 6 \& 8]{douenne_yellow_2022}.


\section{Willingness to contribute: substantial support when citizens bear part of the cost}

%===============================================================================================
% - we are not interested in WTP because what matters is willingness to support policy (which makes the richest pay, not everyone for fairness reasons) and willingness to contribute in case it's a policy and the others are also put to contribution
%===============================================================================================

This section examines which international policies citizens are willing to support and how their support varies with policy design, including contribution rules, country participation, effectiveness, fairness, and redistribution. Accordingly, the relevant outcome is policy support rather than the amount an individual is willing to pay. \textit{Willingness to pay} captures an individual's trade-off between a monetary sacrifice and the achievement of a given objective, whereas \textit{willingness to support} refers more broadly to the acceptance of a collective policy. Individuals may support a policy without being net contributors. For many of the international policies examined thus far, contribution rules protect lower-income households and assign a larger share of the costs to higher-income individuals, wealthier taxpayers, and countries with greater financial capacity. Public support therefore depends partly on considerations of fairness, especially the progressivity of contributions and the distribution of the fiscal burden within and across countries. \textit{Willingness to support} is thus the more appropriate concept for international policies, since contributions vary with ability to pay and cannot be represented by a uniform \textit{willingness to pay}. The remainder of this section therefore examines how far citizens are prepared to support such collective arrangements, moving from a uniform contribution of 1\% of income to the use of domestic taxation for global problems and, finally, to a self-designed global redistribution that may reduce respondents' own income.

%===============================================================================================
% - 1% for climate action (Andre et al)
%===============================================================================================

%LB: J'ai aussi pris en compte les commentaires que tu m'avais laissé dans la V1 du draft.
Drawing on the 2021--2022 Gallup World Poll, \citet{andre_globally_2024} analyze responses from representative samples of people aged 15 and older across 125 countries. The survey asks respondents whether they would contribute 1\% of their household income each month to fight global warming. In the population-weighted sample, 69\% accept this contribution. An additional 6\% would contribute a smaller amount, while 26\% would contribute nothing. Willingness to contribute the full 1\% exceeds 50\% in 114 of the 125 countries and ranges from 30\% to 93\%. The eleven countries below the majority threshold are Japan (49\%), Canada (49\%), the United States (48\%), the United Kingdom (48\%), Pakistan (47\%), New Zealand (46\%), Kazakhstan (45\%), Russia (41\%), Lithuania (41\%), Israel (37\%), and Egypt (30\%) \citep[Figure 1 \& Supplementary Table S4]{andre_globally_2024}. 

Misperceptions about collective participation may help explain some of the reluctance to contribute. Although 69\% of respondents are willing to contribute, they estimate that only 43\% of their fellow citizens would do so, a gap of 26 percentage points. Actual willingness exceeds perceived willingness in every surveyed country. Moreover, respondents' beliefs about others' willingness correlate strongly with their own willingness to contribute \citep[Figure 3 \& Supplementary Table S11]{andre_globally_2024}.

%/LB: p.10-11 du rapport de la WB (2009) => argument global governance (section 5, Ghassim) via 'Evaluation Government Action'.
An earlier World Bank survey covering 15 countries yields similar results. Respondents first indicated whether they would accept higher prices for energy and other products to fight climate change. Within each country, all respondents were presented with the same nominal monthly amount, calculated as 1\% of that country's annual GDP per capita and prorated monthly; those who rejected it then considered a second fixed amount equivalent to 0.5\% of GDP per capita. Because the amount did not vary with individual income, this design is regressive within countries and may therefore yield conservative estimates of support relative to a proportional contribution. Willingness to accept the initial amount ranged from 11\% in Russia to 68\% in China. Acceptance of the full amount was highest in China (68\%), Vietnam (59\%), Japan (53\%), and Iran and Mexico (51\% each), whereas it was lowest in Russia (11\%) and Bangladesh (32\%). The pattern thus does not follow a simple divide between richer and poorer countries. When the lower amount is included, support reaches a majority in 14 of the 15 countries and is highest in Vietnam (85\%) and China (82\%), while Russia is the only country without majority support (25\%). Overall, 46\% accepted the full amount, 63\% accepted either amount, and 33\% rejected both \citep{world_bank_public_2009}. Because all respondents within a country were asked to pay the same absolute amount regardless of their income, the proposal is regressive, imposing a proportionally greater burden on poorer individuals and thus probably yielding a conservative estimate of support.
 % AF: so they propose the same absolute amount to everyone instead of a relative (%) amount? If so, we should flag this as probably yielding conservative results given the regressivity of the proposal. Also, give more info on the country specific attitudes (which types of country have low vs. high support)
%LB: respondents face the same relative amount (1\% of their country GDP per capita), so there's indeed an issue of regressivity.

%===============================================================================================
% - Most agree own taxes should pay for global issues (Fabre 26)
%===============================================================================================

%LB: j'ai ajouté un contraste avec ce qu'on avait dit avant dans la section HICs financing LICs. Ça fait un rappel avec le premier paragraphe de cette section.
\citet{fabre_public_2026} asked respondents whether they agreed that their ``taxes should go towards solving global problems''. Excluding neutral responses, 59\% agree and 41\% disagree \citep[Figures S59--S60]{fabre_public_2026}.
Earlier results show substantially higher acceptance when the contribution falls on the richest: policies requiring contributions from the global top 1\% and top 3\% receive acceptance rates of 69\% and 64\%, respectively. This contrast suggests that support depends on how the policy distributes the contribution burden.

%===============================================================================================
% - Almost half choose a custom redistr that makes them lose (Fabre 26)
%===============================================================================================

%LB: ATTENTION c'est quelque chose que j'ai déjà cité dans la section 4, donc il faudra l'enlever et la remettre ici (j'imagine). J'ai fait un copier-coller de la section 4.
\citet{fabre_public_2026} then asked respondents to construct their preferred global distribution of after-tax income by choosing the shares of winners and losers and the redistribution from richer to poorer individuals. Respondents see both the current and proposed distributions: ``The current distribution is in red, and your custom one is in green'' and could preserve the status quo. The interactive exercise is available at \url{http://preferences-pol.fr/custom_global_redistr.html}, with an explanatory video at \url{https://youtu.be/gSfsQwczT7w}. The median preferred parameters are to make 49\% of the world population better off and 18\% worse off, transferring 5\% of world income from richer to poorer individuals and raising the minimum monthly income to \$287 at purchasing power parity \citep[Figure S49]{fabre_public_2026}. Almost half of respondents (48\%) select a redistribution that would reduce their own income, compared with 9\% who would personally gain. Only 10\% choose to retain the current distribution and not reduce inequality. The average preferred redistribution would likewise transfer 5.4\% of world income, in this case from the top 27\% to the bottom 73\%, and establish a minimum monthly income of \$247 \citep[Figure S50]{fabre_public_2026}. The average and median proposals would reduce the global Gini coefficient from 0.67 to 0.59 and 0.58, respectively. 

\section{Explaining the lack of prominence: low salience, pluralistic ignorance, and national framing}
%===============================================================================================
% - not a lack of sincerity (petition, list experiment, conjoint analyses)
%===============================================================================================

%LB: ici le paragraphe introductif manquant en reprenant tes commentaires de la version passée. 
The evidence reviewed above gives rise to a preference--prominence puzzle: globally redistributive policies attract majority acceptance when surveys present them to citizens, yet they remain largely absent from public debate and political programmes. \citet{fabre_majority_2025} and \citet{fabre_public_2026} examine several explanations for this discrepancy. Petition and list-experiment results assess whether insincerity or social-desirability bias inflates responses to direct survey questions. Conjoint analyses test whether support persists when respondents choose among competing policy bundles or candidates. Evidence from the Green New Deal illustrates how an adversarial political campaign can erode support. Measures of second-order beliefs identify pluralistic ignorance, while responses to unprompted open-ended questions reveal the low salience of global redistribution. Finally, an election-framing experiment tests whether political institutions organized at the national level direct citizens' attention toward national constituencies and policy responses. Taken together, these results provide little evidence of widespread preference falsification. Instead, they point to three complementary explanations for the limited political prominence of international redistribution: pluralistic ignorance, low salience, and national political framing.

%LB: Petition - Fabre 25
In a petition experiment, \citet{fabre_majority_2025} find that willingness to sign a petition supporting the GCS is nearly as high as support in the case where respondents are asked directly, and that the gap between the two methods is similar for National Redistribution (NR)\footnote{NR: National Redistribution policy.}. In the United States, 51\% would sign the GCS petition, compared with 53\% who support it in the direct question. In Europe, the corresponding figures are 69\% and 76\%. The differences are comparable for NR: 57\% would sign the petition and 58\% support it directly in the United States, while the corresponding figures are 67\% and 72\% in Europe. Spanish respondents report the highest willingness to sign, at 78\% for the GCS and 74\% for NR \citep[Figure S27]{fabre_majority_2025}.

%LB: List experiment - Fabre 25
\citet{fabre_majority_2025} also conduct a list experiment to assess whether social desirability inflates stated support for the GCS\footnote{\textit{Social desirability bias} occurs when respondents report opinions that present them in a favorable light.}. Respondents report only how many policies they support, while the researchers randomly vary the composition of the list and by including or not the GCS. This design estimates tacit support for the scheme without requiring respondents to endorse it explicitly. The experiment yields an estimated support rate of 62\% overall, including 52\% in the United States and 72\% in Europe. Since these estimates are not significantly different than direct stated support, the authors find no evidence of social desirability bias \citep[Table 1]{fabre_majority_2025_si}. 

%LB: conjoint analysis - Fabre 25 - version étendue 
\citet{fabre_majority_2025} then use a series of conjoint analyses in which respondents choose between policy bundles or candidates' platforms. In the first analysis, respondents choose between a bundle combining the GCS, NR and a national climate policy: (C) (a coal phase-out in the U.S.) and an otherwise identical bundle containing only NR and (C). The bundle including the GCS is preferred by 55\% of U.S. respondents and 74\% of European respondents. 
In a second analysis, respondents are randomly assigned to comparisons in which NR is held constant while the bundles add the GCS, (C), or both. Support for the GCS conditional on NR reaches 55\% in the United States and 77\% in Europe and is not significantly different from support for the GCS alone. Moreover, when respondents choose directly between NR combined with the GCS and NR combined with (C), 53\% of U.S. respondents prefer (C), while 52\% of European respondents prefer the GCS, neither difference is statistically significant. These results indicate that support for the GCS depends neither on compensation through NR nor on its combination with a national climate policy \citep[Figure S8]{fabre_majority_2025}. 
A third analysis presents respondents with progressive and conservative candidates' platforms and randomizes only whether the progressive platform includes the GCS. Adding the GCS increases the probability of being chosen in any country by 2.8 percentage points on average and by 11.2 points in France \citep[Supplementary Table 2]{fabre_majority_2025}.

%===============================================================================================
% - campaign effect (Gustafson et al)
%===============================================================================================

Drawing on two surveys of U.S. voters conducted in December 2018 and April 2019, \citet{gustafson_development_2019} examine how attitudes toward the Green New Deal (GND) changed after the proposal entered the national political debate. Each survey presented respondents with a description of the proposal and asked: ``How much do you support or oppose this idea?'' In December, 81\% supported the GND, including 92\% of Democrats and 64\% of Republicans. By April, the share of voters who had heard at least ``a little'' about the proposal had risen from 17\% to 59\% \citep[Figure 1]{gustafson_development_2019}. Over the same period, support declined among Republicans but remained stable among Democrats. It fell from 64\% to 44\% among Republicans \citep[Figure 2]{gustafson_development_2019}. 
In April, Republican support varied sharply with prior exposure: 85\% of those who had heard nothing about the GND supported it, compared with only 4\% of those who had heard a lot \citep[Figure 3]{gustafson_development_2019}. Among frequent Republican viewers of Fox News, support fell from 54\% to 22\%, compared with a decline from 71\% to 56\% among infrequent viewers. 

%===============================================================================================
% - pluralistic ignorance
%===============================================================================================

Furthermore, \citet{fabre_public_2026} ask respondents whether they accept the Global Climate Scheme (GCS) and then elicit their second-order beliefs about the share of their compatriots who accept it. The median respondent underestimates national acceptance by 16 percentage points, while non-American respondents underestimate acceptance in the United States by 22 points \citep[Figure S30]{fabre_public_2026}. These beliefs correlate strongly with respondents' own attitudes: 51\% of GCS supporters believe that a majority of their compatriots also accept the scheme, compared with only 24\% of opponents. The author reports that ``support for the GCS reaches 72\% among the 39\% of respondents who believe that a majority in their country supports it, compared to 44\% among those who do not'' \citep[Figure S30]{fabre_public_2026}. These associations may reflect both pluralistic ignorance and a ``false consensus effect''.

%===============================================================================================
% - low salience => we should talk about "acceptance" rather than "support"
%===============================================================================================

Before the survey revealed its topic, respondents were randomly assigned one of four broad open-ended questions asking about their main concerns, their needs or wishes, an issue they considered important but neglected in public debate, or the greatest injustice. Among responses to the questions about concerns and wishes, 31\% mentioned ``money'' or the ``cost of living''. Among responses to the greatest-injustice question, only 3.2\% explicitly referred to global inequality, while global redistribution almost never appeared as a wish \citep[Figure 3]{fabre_public_2026}. Taken together, these findings reveal broad acceptance when respondents evaluate a specific proposal but low salience in their spontaneous political concerns. They do not establish that this acceptance translates into active support capable of shaping votes or political mobilization.


%===============================================================================================
%- national framing: original results (https://github.com/bixiou/budget_26/blob/main/figures/group_defended.pdf
%Weighted regression:
%                           Estimate Std. Error t value Pr(>|t|)   
%(Intercept)                  0.16476    0.05966   2.762  0.00586 **
%variant_group_defendedworld  0.16152    0.08268   1.954  0.05103 . 
%---
%Signif. codes:  0 ‘***’ 0.001 ‘**’ 0.01 ‘*’ 0.05 ‘.’ 0.1 ‘ ’ 1
%
%Residual standard error: 1.306 on 998 degrees of freedom
%Multiple R-squared:  0.003809,	Adjusted R-squared:  0.002811 

%Non-weighted:
%                                               Estimate Std. Error t value Pr(>|t|)   
%(Intercept)                                0.15933    0.05952   2.677  0.00755 **
%variant_group_defendedworld  0.20014    0.08230   2.432  0.01520 * 
%---
%Signif. codes:  0 ‘***’ 0.001 ‘**’ 0.01 ‘*’ 0.05 ‘.’ 0.1 ‘ ’ 1

%Residual standard error: 1.3 on 998 degrees of freedom
%Multiple R-squared:  0.005891,	Adjusted R-squared:  0.004895 
%)
%===============================================================================================

A final explanation is a nationalist bias. Fabre in an original survey experiment provides evidence for this mechanism. Respondents were randomly asked either which group they defend when voting or which group they would defend in hypothetical world elections. Responses were coded on a scale ranging from (-2) for ``Family and self'', through \(0\) for ``Fellow citizens'', to (2) for ``Sentient beings''. Framing the elections at the world level shifted responses by 0.162 points toward more universalist groups in the weighted regression, both estimates are significant (\url{https://github.com/bixiou/budget_26/blob/main/figures/group_defended.pdf}). 

\newpage
\bibliographystyle{plainnat}
\bibliography{chapter/references}

\end{document}
